\subsection{A sealed current-STA resource trajectory and its costs}
The independent pooled-analysis controller completes two accepted transactions from the automatically obtained R1 fast state, with final context/input seals and complete native/physical/STA gates. Exact CPD stays13.228145689936355 ns. Final wirelength is2068988, with85 SB and5071 CB resources:108 more wire units, one fewer SB and one fewer CB than the fast root. The observed joint quality population now contains198 records. The new states quality-dominate the earlier R1 sharing point (13.421610 ns,2069004 wire,86 SB,5071 CB) and historical N (13.4216 ns,2069228 wire,90 SB,5071 CB), accounting for printed-CPD intervals. Strong native A0/R1 remain in this joint population. The quality-frontier group count can stay constant while two new groups replace two inferior groups; count alone is not evidence of no progress. These are observation-level comparisons, not a history-best configuration policy or matched complete-runtime superiority.

\begin{table}[t]
\centering\scriptsize
\caption{Own-current R1 resource trajectory. All listed candidate routes are legal; exact CPD is unchanged.}
\begin{tabular}{lrrrr}
\hline
State/arm & Wire & SB & CB & Commit\\
\hline
Fast root &2068880&86&5072&initial\\
First OFF &2068880&86&5072&no\\
First ON &2068976&86&5071&yes\\
Second OFF &2068988&85&5071&control\\
Second ON &2068988&85&5071&yes\\
\hline
\end{tabular}
\end{table}
The selector first chooses net31972/pin8. The first unrestricted native transaction changes no resources; the crossing-domain candidate removes one CB at96 extra wire units. It regresses one setup component, still4.869294 ns above the new worst slack. After committing, complete current STA is reused to recompute the rank: eligible scopes change from62 to60, and the next chosen target is net21922/pin4. Both OFF and ON remove one SB at twelve additional wire units, with unchanged endpoint setup/hold vectors. The ON final candidate reuses the exact OFF native analysis observation. Only the registered resource-envelope mode accepts these moves; strict CPD sees no improvement. The first move fails the endpoint-vector rule, and the second leaves the endpoint vector flat. Near-endpoint counts within0.05/0.1/0.2/0.5/1.0 ns remain4/4/5/10/77. The selected worst endpoints, complete winning arrival cones, and ten winning routed-connection identities are unchanged after both commits, so there is no observed critical-path migration on this trajectory.

Attribution must follow the controls. The first local comparison supports the domain restriction, whereas the second supports explicit current-state selection using the existing native kernel. It does not establish that the first release enabled the second move: the initial-state native-only trajectory is a separate queued ablation with its own state and successor targets. The method's candidate-domain/selection/scheduling contribution is not a new connection-search primitive. Other legal rejected tradeoffs remain archived.

The complete imported-state paired-control workflow takes1363.975 s. Measured sequential phase sums are556.965 s for native wrappers,148.520 s for transaction audits,21.850 s for domain audits,424.472 s for pooled candidate FULL/STA/materialization, and55.621 s for current selection. The residual covers initial preparation/observation and final sealing. The four native analysis children total152.807 s and native repair cores total3.528 s; these are nested components, not additional costs to sum with the workflow. One native observation is exactly reused. The original route, cold complete single-configuration algorithm time, and a formally rotated timing comparison are still outside this measurement. This cost profile motivates both fewer repeated state constructions and transaction batching; neither is already a demonstrated speedup.

\paragraph{Independent native-worker reuse.}
A new201.620 s build extends the pre-route private queue with per-child native transaction requests, objective, crossing domain and pressure. Those hooks are set only after fork; the parent remains free of route and native-repair state. The native resource-domain kernel is unchanged. An existing two-layer fixture passes same-binary standalone versus worker checks for unrestricted, retained-crossing and infeasible transactions, and retains full-precision original STA after an invalid native job. The check executes eight mixed worker jobs and three exact cache hits in16.988 s, with no new packing, placement or routing flow. Mutable request contents have separate keys, and deliberate corruption of a separate test-owned cached output is rejected and fails its final seal. A worker's STA describes its input route; private emitted trees still need actual replacement, independent capacity/scope audits and a separate complete native FULL/STA before acceptance. Large BNN worker/replacement equivalence is queued after the current quality owner. The following independent controller and K4/K8/native-only cases are already registered behind that validation. Build/tiny-test success does not establish BNN equivalence, algorithm quality or complete runtime improvement.

\paragraph{Transfer opportunity before more routing experiments.}
The same read-only current selector is applied to matched A0/R1/Tf physical inputs and native constraints for two additional circuits. Exact existing resource counts are reproduced. For spmv, strong A0 has no multi-crossing net and no eligible sharing scope; R1 and Tf each have two scopes. For eltwise, A0/R1/Tf have12/21/75 scopes. These counts are screening opportunities, not achievable savings. Inactivity on strong spmv A0 limits this sharing interface, not the whole wirelength/timing research space. The new development cohort includes strong eltwise A0 and spmv R1 controls, alongside the stronger BNN A0 start, and never chooses a different history-best configuration for each circuit.

A proposed distinct-net bundle takes the first ranked cut from each selected net, avoiding same-net cut overlap. Native restoration updates real occupancy after every connection; the complete combination receives one final budget/FULL/STA decision and either commits or rolls back. K1 requests reproduce the current single-cut request exactly on the registered source; K4/K8 request generation and rank windows are verified but their feasible savings are unknown. This is an intended intervention on repeated full-validation cost, not a reason to skip final STA or require every component to improve CPD independently before a combination is considered.
